\subsection{从辐射通道到有限孔径计数}
\label{chap:feqo-cl}

\begin{proposition}[偶极轴向圆锥的收集比例]
设真空远场由固定方向的电偶极 $\mathbf p=p_z\hat{\mathbf z}$ 辐射，探测器完整收集以 $+z$ 为轴、半角 $\theta_{\max}$ 的圆锥内两种偏振。其在总远场辐射中的比例为
\begin{equation}
\eta_{\rm cone}
=\frac{2-3\cos\theta_{\max}+\cos^3\theta_{\max}}4.
\label{eq:feqo-dipole-collection-fraction}
\end{equation}
\end{proposition}
\begin{proof}
偶极远场系数为
\[
\mathbf F(\hat{\mathbf n},\omega)
=\frac{k^2}{4\pi\varepsilon_0}
\hat{\mathbf n}\times(\mathbf p\times\hat{\mathbf n}),
\qquad |\mathbf F|^2\propto\sin^2\theta.
\]
因此角度功率积分给
\[
\eta_{\rm cone}
=\frac{2\pi\int_0^{\theta_{\max}}\sin^3\theta\dd\theta}
{2\pi\int_0^\pi\sin^3\theta\dd\theta},
\qquad\int_0^\pi\sin^3\theta\dd\theta=\frac43,
\]
积分后即得结论。半球给 $1/2$；窄轴向圆锥满足 $\eta_{\rm cone}=3\theta_{\max}^4/16+O(\theta_{\max}^6)$，体现偶极轴上的辐射零点。
\end{proof}

\begin{workedexample}[辐射分支与孔径效率]
若电子交出的能量中 $80\%$ 成为远场光，而探测器收集其中 $25\%$，有效收集比例为 $0.80\times0.25=0.20$。在上述轴向偶极模型中，$\theta_{\max}=60^\circ$ 给 $\eta_{\rm cone}=5/32$；若另有 $80\%$ 辐射分支，则总损失中被收集的比例为 $1/8$。

对于有导波出口的结构，稳态能量平衡应写成
\[
W_{\rm loss}=W_{\rm abs}+W_{\rm far}+W_{\rm guided},
\qquad W_{\rm det}\le W_{\rm far}.
\]
仅在无其他出口、初末局域储能相同的条件下，才能简化为吸收加远场辐射两项。数值孔径改变收集投影；材料吸收和导波损失还需由各自通道求出。
\end{workedexample}
